
%! Author = stanislaw
%! Date = 12.03.2026

% Preamble
\documentclass[11pt]{article}

% Packages
\usepackage{amsmath}
\usepackage{tabularx}
\usepackage{booktabs}
\usepackage{seqsplit}

% Document
\begin{document}

\begin{table}[h!]
    \centering
    \small
    \begin{tabularx}{\textwidth}{|l|X|X|}
        \hline
        File name & \seqsplit{aut\_branching} & \seqsplit{basic\_branching} \\
        \hline
		cubes3.grl & 0.13 & 0.07 \\
		cubes4.grl & 0.74 & 0.77 \\
		cubes5.grl & 5.16 & 18.2 \\
		trees36.grl & 0.27 & 0.18 \\
		wheeljoin14.grl & 2.44 & 7.62 \\
		modulesD.grl & 183.33 & 409.57 \\
        \hline
    \end{tabularx}
    \caption{
        This run compares branching using generator-based automorphism counting v1\_0\_5 against basic branching v1\_0\_3 with naive automorphism counting.
        The generator-based automorphism counting does add some overhead, but it significantly reduces the computation time for large \& symmetric graphs
    }
    \label{tab:1}
\end{table}

\end{document}
